\documentclass[11pt]{article}
\usepackage{amsmath,amssymb,amsthm}
\usepackage[margin=1in]{geometry}
\usepackage{booktabs,longtable,array}
\usepackage[colorlinks=true,linkcolor=blue,urlcolor=blue]{hyperref}

\newtheorem{theorem}{Theorem}
\newtheorem{lemma}[theorem]{Lemma}
\newtheorem{proposition}[theorem]{Proposition}
\theoremstyle{definition}
\newtheorem{definition}[theorem]{Definition}
\theoremstyle{remark}
\newtheorem{remark}[theorem]{Remark}

\newcommand{\eps}{\epsilon}
\newcommand{\nh}{\hat n}
\newcommand{\dd}{\mathrm{d}}
\newcommand{\sky}[1]{\langle #1\rangle_{\Omega}}
\newcommand{\ens}[1]{\langle #1\rangle}
\newcommand{\T}{\mathcal{T}}
\newcommand{\TF}{\operatorname{TF}}
\newcommand{\Cin}{\operatorname{Cin}}
\newcommand{\Si}{\operatorname{Si}}
\newcommand{\Ephi}{\mathfrak{e}}
\newcommand{\st}[1]{\textsf{[#1]}}
\newcommand{\code}[1]{\texttt{#1}}

\title{The observed temperature power spectrum through fourth order\\
in a radiation-only universe: final report}
\author{Unified analysis (branch \code{unified/temperature-spectrum})}
\date{2026-10-31}

\begin{document}
\maketitle

\begin{abstract}
We compute the observed temperature anisotropy through third order, and the angular power
spectrum through fourth order, in the primordial amplitude~$\eps$. The setting is a flat
radiation-only universe in general relativity, with Planck test photons emitted on
$T_{\rm phys}=T_\star$ and an observer on a fluid worldline.

The calculation uses the fluid's velocity potential and comoving labels as coordinates
(\emph{temperature frame}). In these coordinates the observed temperature is an exact line
integral of the shear of a conformally rescaled metric, and the observer's own temperature
cancels exactly in the normalised map. Further results:
\begin{itemize}
\item The spatial Einstein equations take a new mixed trace-free form, which yields a
  triangular solution scheme with elementary Green functions at every order.
\item This scheme satisfies all ten field equations at every order (proved).
\item All second-order field kernels are given in closed form for general momenta, in the
  entire functions $\Si$, $\Cin$ and $\Ephi_m$. This makes every resonance surface manifestly
  regular.
\item Third-order kernels are explicit nested integrals.
\item The complete map $\Theta^{(1,2,3)}$ is constructed. It retains photon propagation,
  emission surface, observer clock, tetrad and normalisation terms.
\item $C_\ell^{(11)}$, $C_\ell^{(22)}$ and $C_\ell^{(13)}$ are reduced to exact
  finite-dimensional integrals, with every sky integral done in closed form.
\item In the superhorizon limit the complete observable is obtained in closed form,
  \emph{exact in~$\zeta$}. It gives
  \[
  C_2=\tfrac{4\pi}{225}\eps^2\sigma^2\Xi^2\big[1+\tfrac{29}{3}\eps^2\sigma^2\big]
  +O(\eps^6,k_0^6),\qquad \Xi\equiv x_o^2-x_\star^2 .
  \]
\end{itemize}
Every displayed identity is checked by exact computer algebra (no numerics). We state
explicitly which integrals remain unevaluated, and which limits remain open.
\end{abstract}

\tableofcontents

%======================================================================
\section{Problem and summary of results}
%======================================================================

The task (\code{problem\_statement.tex}) asks for $C_\ell=\eps^2C_\ell^{(11)}+\eps^4[\frac14
C_\ell^{(22)}+\frac13\operatorname{Re}C_\ell^{(13)}]+O(\eps^6)$ for the observed map
\[
\Theta(\nh)=\frac{T_{\rm obs}(\nh)}{\frac1{4\pi}\int \dd\Omega'\,T_{\rm obs}(\nh')}-1,\qquad
T_{\rm obs}=T_\star\frac{E_o}{E_e},\quad E=-u_\mu p^\mu .
\]
The physical model is a perfect radiation fluid with $P=\rho/3$, $\rho=\alpha T_{\rm phys}^4$
and $\Lambda=0$, with Gaussian adiabatic primordial perturbations and no independent
higher-order modes. The calculation must be analytical and keep every term, with all
normalisations fixed. A monochromatic specialisation is to be made only at the end.

\begin{center}
\begin{longtable}{p{6.4cm}p{8.6cm}}
\toprule
Requirement & Result (section; status) \\\midrule\endhead
Gaussian variable; no higher-order modes & $\zeta=\frac16\ln\det\gamma_{ij}$, exactly conserved; Def.~\ref{def:ic} (\S\ref{sec:ic}) \\
Constrained initial data at every order & regular solution at any $t_i$; constraints proved (Thm.~\ref{thm:consistency}); forced data nonzero \\
Field solution through $\eps^3$ incl.\ vectors/tensors & triangular scheme (\S\ref{sec:fields}); 2nd order \st{closed form}; 3rd order \st{explicit nested integrals} \\
Observer, clock, tetrad, emission surface & \S\ref{sec:obs}; the tetrad reduces to a symmetric root (Prop.~\ref{prop:frame}) \\
Complete $\Theta^{(1)},\Theta^{(2)},\Theta^{(3)}$ & \S\ref{sec:obs}, with an independent generator and Ward checks \\
$C_\ell^{(11)}$ & \st{closed form} ($\ell=2,3$; general $\ell$ as a 1D integral) \\
$C_\ell^{(22)}$, $C_\ell^{(13)}$ & \st{exact finite-dimensional integrals}, sky integrals evaluated (\S\ref{sec:spec}); existence proved \\
Monochromatic, held-fixed quantity & $\ens{g^2}=\sigma^2$ fixed (\S\ref{sec:mono}) \\
Asymptotic limits & superhorizon: \st{closed form, exact in $\zeta$}; late time: \st{first order}; general $k_0$ at fourth order: \st{open} (\S\ref{sec:limits}) \\
Convergence for general spectra & compact support: \st{proved}; IR: leading cancellation \st{derived} from the Ward identities; UV: \st{power counting only} \\
Verification & \S\ref{sec:ledger} (exact symbolic, generic-jet proofs) \\
\bottomrule
\end{longtable}
\end{center}

%======================================================================
\section{Conventions and the experiment}\label{sec:conv}
%======================================================================

Signature $(-,+,+,+)$, $c=1$. The background is $a=a_1t$ and $\bar T=\T/(a_1t)$, with
$8\pi G\alpha\T^4=3a_1^2$.

\paragraph{Fourier and perturbative conventions.}
\begin{itemize}
\item $g(X)=\int_k g_k e^{ik\cdot X}$, with $\int_k=\int\dd^3k/(2\pi)^3$ and
  $\ens{g_kg_{k'}}=(2\pi)^3\delta^3(k+k')P(k)$; also $\mathcal P(k)=k^3P/2\pi^2$.
\item Every field is $Y=\sum_r\eps^rY^{(r)}/r!$ with
  $Y^{(r)}=\int_{k_1\ldots k_r}K_r(k_1,\ldots,k_r)\,g_{k_1}\cdots g_{k_r}e^{iK\cdot X}$ and
  $K=\sum k_a$. The symmetric kernels carry no extra $1/r!$.
\item \emph{Multilinear rule:} $K_r$ is the coefficient of $z_1\cdots z_r$ when
  $g=\sum_az_ae^{ik_a\cdot X}$ is inserted and only monomials in distinct labels are kept. This
  fixes every combinatorial factor.
\item Spherical harmonics $Y_{\ell m}$ are complex, with the Condon--Shortley phase.
\end{itemize}

\paragraph{Measurement labels.}
\begin{itemize}
\item \emph{Emission:} the surface $T_{\rm phys}=T_\star$. Emission is isotropic and Planckian
  in the fluid rest frame. The intersection is the branch continuously connected to the
  background one.
\item \emph{Observer:} the fluid line $X=0$. The initial event is $T_{\rm phys}=T_I>T_\star$ on
  that line, and observation follows after fixed proper time~$\tau_o$.
\item \emph{Tetrad:} Fermi--Walker transported, with initial orientation
  $\gamma(t_I,0)^{-1/2}$ (the symmetric square root, aligned with the marked labels).
\item \emph{Background times:} $\bar t_\star=\T/(a_1T_\star)$, $\bar t_I=\T/(a_1T_I)$ and
  $\bar t_o^2=\bar t_I^2+2\tau_o/a_1$. We write $x\equiv kt/\sqrt3=k/(\sqrt3\,aH)$.
\end{itemize}

%======================================================================
\section{Exact framework}\label{sec:exact}
%======================================================================

\begin{lemma}[Temperature frame]\label{lem:frame}
For an irrotational radiation fluid, $Tu_\mu=-\T\partial_\mu t$ defines the velocity potential
$t$. In coordinates $(t,X^i)$, with comoving labels $X^i$,
\[
g=N^2\big(-\dd t^2+\gamma_{ij}\dd X^i\dd X^j\big),\qquad N=\T/T_{\rm phys},\qquad u=N^{-1}\partial_t
\]
exactly. Euler's equation holds identically, and continuity is equivalent to
$\partial_t\det\gamma=0$. Hence $\gamma=e^{2\zeta(X)}\exp H$ with $\operatorname{tr}H=0$, and
the shear $\sigma_{ij}=\frac12\partial_t\gamma_{ij}$ is traceless.
\end{lemma}

\begin{lemma}[Exact redshift]\label{lem:redshift}
Null geodesics of $g$ and of $\tilde g=g/N^2$ coincide, and $E=\tilde E/N$. Hence
\[
T_{\rm obs}(\nh)=T_{\rm fluid}(o)\,e^{-I(\nh)},\qquad
I=\int_{t_e(\nh)}^{t_o}\sigma_{ij}(t,X(t))\,e^ie^j\,\dd t,
\]
where $\dd X^i/\dd t=e^i$, $\gamma_{ij}e^ie^j=1$ and
$\dd e^i/\dd t=-2S^i{}_je^j-\Gamma^i{}_{jk}e^je^k+e^i\sigma_{jk}e^je^k$, with
$S=\gamma^{-1}\sigma$. The emission event is the root of $\ln t_e+\nu(t_e,X(t_e))=\ln\bar t_\star$,
with $N=a_1te^\nu$.
\end{lemma}

$T_{\rm fluid}(o)$ is direction-independent. It therefore cancels from $\Theta$ \emph{exactly}.
The emission-surface temperature shift, Sachs--Wolfe, Doppler and integrated Sachs--Wolfe terms
are all contained in the single integral~$I$.

\begin{proposition}[Unified field equations]\label{prop:fe}
With $n=\ln N$, Einstein's equations are equivalent to
\begin{align}
&(N)\quad 3\nu''-\Delta_\gamma\nu-2|D\nu|^2+\operatorname{tr}S^2+\tfrac{3}{t^2}(e^{-2\nu}-1)=0,\\
&(T)\quad \TF\big[S'+2(t^{-1}+\nu')S+R^k{}_j[\gamma]-2D^kD_j\nu+2D^k\nu D_j\nu\big]=0,\\
&(M)\quad D_jS^j{}_i-2\partial_i\nu'+\tfrac2t\partial_i\nu+2S^j{}_i\partial_j\nu+2\nu'\partial_i\nu=0,\\
&(\mathrm{Ham})\quad R-4\Delta\nu-2|D\nu|^2-\operatorname{tr}S^2+6(t^{-1}+\nu')^2-6e^{-2\nu}/t^2=0.
\end{align}
In the mixed form (T), the $\sigma\sigma$ terms of the spatial equation cancel, because
$\gamma^{ki}\sigma'_{ij}=(S^k{}_j)'+2S^k{}_iS^i{}_j$. This is proved for generic
$\gamma$ (six free functions, $t$-independent determinant) and arbitrary $N$, on independent
symbolic 2-jets.
\end{proposition}

%======================================================================
\section{Primordial state and initial data}\label{sec:ic}
%======================================================================

\begin{definition}[Gaussian variable; absence of independent higher-order modes]\label{def:ic}
The state is the solution of Prop.~\ref{prop:fe} with the following properties:
\begin{itemize}
\item[(G)] $\zeta=\eps g$ exactly, with $g$ a zero-mean, homogeneous, isotropic Gaussian field;
\item[(T)] $\gamma\to e^{2\zeta}\delta_{ij}$ as $t\to0$. The limit is conformally flat in the
  marked conformally Cartesian material chart, so there is no tensor seed;
\item[(V)] $\dd(Tu)=0$, so there is no vorticity seed;
\item[(R)] $\nu$ and $H$ are even and analytic in $t$ at fixed momenta, with $\nu(0)=0$ and
  $H(0)=0$, so there are no decaying modes.
\end{itemize}
``No independent higher-order modes'' means $\zeta^{(r\ge2)}\equiv0$ together with (T), (V) and
(R) at every order.
\end{definition}

Because $\det\gamma$ is $t$-independent, $\zeta=\frac16\ln\det\gamma$ at \emph{every} time. As
$t\to0$, constant-$t$ slices become uniform-density slices, so $\zeta$ is the nonlinear
uniform-density ($\delta N$) curvature perturbation.

The residual conformal relabellings are fixed as follows: the dilation by $\ens{g}=0$; special
conformal transformations by statistical homogeneity; translations and rotations are
statistical symmetries.

Exactly in $\zeta$, the gradient-expansion data are
\begin{equation}\label{eq:grad}
\nu=\tfrac{t^2}{18}e^{-2\zeta}\big(2\nabla^2\zeta+|\nabla\zeta|^2\big)+O(t^4),\quad
H_{ij}=\tfrac{t^2}{3}e^{-2\zeta}\,\TF\!\big[\partial_i\partial_j\zeta-\partial_i\zeta\partial_j\zeta\big]+O(t^4).
\end{equation}
The initial data at any $t_i$ are the regular solution of \S\ref{sec:fields} evaluated at
$t_i$, with $\rho=3a_1^2/(8\pi GN^4)$ and $u=N^{-1}\partial_t$. Higher-order data are nonzero.
For example, $\gamma^{(2)}(0)=4g^2\delta$, and the $\nu^{(2)}$ kernel is
$\frac{t^2}{9}[2(k_1^2+k_2^2)-k_1\cdot k_2]+O(t^4)$. The forced $K=0$ mode ($5k^2t^2/9$ at
$k_2=-k_1$) is kept. The data evolve back into the same regular solution, so $t_i$ drops out
identically.

%======================================================================
\section{Field solution at every order}\label{sec:fields}
%======================================================================

\subsection{Triangular scheme}

At total momentum $K$, the order-$r$ kernels $(\nu_r,A_r)$, with $A_r$ the order-$r$ part of
$H$, solve $\mathcal L_K[\nu_r,A_r]+c^{(r)}=0$. The operators are the same at every order:
\[
L_N=3\nu''+K^2\nu-\tfrac{6}{t^2}\nu,\quad
L_M=\tfrac i2K_jA'_{ji}-2iK_i(\nu'-\nu/t),\quad
L_T=\tfrac12(A''+\tfrac2tA')+\TF R_1[A]+2\TF(K_iK_j)\nu .
\]
Here $R_1[A]_{ij}=\frac12(K^2A_{ij}-K_aK_iA_{aj}-K_aK_jA_{ai})$. The source $c^{(r)}$ is the
multilinear coefficient with $(\nu_r,A_r)=0$; it is explicit in the lower orders.

Write $c=K/\sqrt3$, $x=ct$, $u_1=\sin x/x-\cos x$ and $u_2=-\cos x/x-\sin x$. The coefficient
$c_r$ is the $t^2$ coefficient of \eqref{eq:grad} at order $r$. In particular
$c_2=[2(k_1^2+k_2^2)-k_1\cdot k_2]/9$ and $c_3=\frac29[-2\sum k_a^2+\sum_{a<b}k_a\cdot k_b]$.
Then
\begin{align}
\nu_r&=\frac{9c_r}{K^2}u_1(x)-\frac{1}{3c}\Big[u_2(x)\!\int_0^t\!u_1(cs)c_N^{(r)}\dd s-u_1(x)\!\int_0^t\!u_2(cs)c_N^{(r)}\dd s\Big],\\
h_S&=\tfrac{3}{2K^2}\big[4K^2\nu_r+12\nu_r'/t+12\nu_r/t^2+c_{\rm Ham}^{(r)}\big]\quad\text{(algebraic, from (Ham))},\\
w&=\tfrac{2i}{K^2}\int_0^tP_\perp c_M^{(r)}\dd s,\qquad
h=-2\int_0^t\tfrac{s}{t}\tfrac{\sin K(t-s)}{K}\,\mathrm{TT}[c_T^{(r)}](s)\,\dd s,\\
A_r&=(\hat K\hat K-\tfrac13\delta)h_S+(K\otimes w+w\otimes K)+h .
\end{align}
At $K=0$ no projector appears: $A_r=-2\int_0^t\frac st(t-s)c_T^{(r)}(s)\dd s$ and
$\nu_r=c_rt^2-\frac13\int_0^t\frac{t^3-s^3}{3ts}c_N^{(r)}\dd s$. The particular parts are
$O(t^4)$, so the free regular lapse mode is fixed \emph{exactly} by (Ham) at $O(t^2)$.

\begin{theorem}[Consistency at every order]\label{thm:consistency}
The triangular solution satisfies all ten field equations at every order.
\end{theorem}
\begin{proof}
Let $\mathbb E_{\mu\nu}=G_{\mu\nu}-8\pi GT_{\mu\nu}$. Matter conservation is exact
(Lemma~\ref{lem:frame}), so $\nabla^\mu\mathbb E_{\mu\nu}=0$. If all residuals vanish below
order~$r$, the order-$r$ residual obeys the identity linearised about $a^2(-\dd t^2+\dd x^2)$,
with $\mathcal H=1/t$. By construction:
\begin{itemize}
\item $\mathbb E^t{}_i=0$, from (M);
\item $\mathbb E^t{}_t=\frac12\mathbb E$, from (N). With $\mathbb E^i{}_i=3\tau$ this gives
  $\mathbb E^t{}_t=3\tau$;
\item for $K\neq0$, the TT part of $\TF\,\mathbb E^i{}_j$ vanishes.
\end{itemize}
Write the remaining trace-free part as $e_S(\hat K\hat K-\delta/3)+\hat K\otimes w+w\otimes\hat K$.
\begin{itemize}
\item The spatial component gives $K_j\mathbb E^j{}_i=0$. Hence $w=0$ and $\frac23e_S+\tau=0$.
\item The time component gives $3\tau'+6\tau/t=0$. Regularity then forces $\tau=0$, so
  $e_S=0$ and $(\mathrm{Ham})\propto\mathbb E^t{}_t=0$.
\end{itemize}
At $K=0$, (T) holds in full, so $\tau=0$. The spatial component gives
$\partial_t\mathbb E^t{}_i+4\mathcal H\mathbb E^t{}_i=0$, so $\mathbb E^t{}_i=0$ by regularity.
\end{proof}

\subsection{Explicit kernels}

\paragraph{First order \st{closed form}.}
\[
\nu_1=-(\sin x/x-\cos x),\qquad
H_1=-\tfrac{6}{k^2}\big(\tfrac{2\sin x}{x}-\cos x-1\big)\big(k_ik_j-\tfrac13\delta_{ij}k^2\big),\qquad
\zeta_1=1 .
\]

\paragraph{Second order \st{closed form, general momenta}.} Every integral above is of the
form $\int_0^tB(s)\dd s$, where $B=\sum_\omega e^{i\omega s}R_\omega(s)$ and each $R_\omega$ is a
Laurent polynomial. The frequencies $\omega$ are integer combinations of $p/\sqrt3$, $q/\sqrt3$,
$K/\sqrt3$ and $K$ (with $p=|k_1|$, $q=|k_2|$, $K=|k_1+k_2|$); they vanish on resonance surfaces.
We integrate term by term in \emph{entire} functions:
\begin{align}
\int_0^t s^me^{i\omega s}\dd s&=t^{m+1}\Ephi_m(i\omega t),\qquad
\Ephi_m(z)=\int_0^1u^me^{zu}\dd u\qquad(m\ge0),\\
\int_0^t s^{-M}\Big[e^{i\omega s}-\sum_{j<M}\tfrac{(i\omega s)^j}{j!}\Big]\dd s&=J_M(\omega,t),\quad
J_1=-\Cin(\omega t)+i\Si(\omega t),\\
J_M&=\tfrac{t^{1-M}E_M}{1-M}+\tfrac{i\omega}{M-1}J_{M-1},\qquad
E_M=e^{i\omega t}-\sum_{j<M}\tfrac{(i\omega t)^j}{j!}.
\end{align}
Here $\Cin(x)=\int_0^x(1-\cos u)u^{-1}\dd u$ is entire and even, and $\Si$ is entire and odd.
The subtracted Taylor pieces sum to the principal part of $B$ at $s=0$, which vanishes (this is
asserted in the code). \emph{No division by any frequency occurs}, so the kernels are manifestly
finite on every resonance surface:
\begin{itemize}
\item the scalar collinear surfaces $\mu=\pm1$;
\item the tensor and vector sound cone $\mu=-1/3$ (for $|k_1|=|k_2|$).
\end{itemize}
The removable $\ln|\alpha|$ terms that $\mathrm{Ci}$ would introduce never appear. The complete
expressions for $\nu_2$, $h_S$, $w$ and $h$ are in
\code{solution/results/second\_order\_closed\_forms.txt}. Each antiderivative is verified by exact
differentiation. Substituting the assembled $\nu_2$ back into the lapse equation (N) gives an
exact zero in closed form. The scalar part $h_S$ satisfies (Ham) by construction.

The momentum and trace-free equations follow from three facts: the verified antiderivatives,
the verified Green functions, and Thm.~\ref{thm:consistency}. They are also confirmed
independently by the exact series test through $O(t^8)$. A direct closed-form substitution into
(M) and (T) is available (\code{--check all}) but was not run to completion: it takes hours, and
it is redundant with these checks.

\paragraph{Third order \st{explicit nested integrals}.} For $C^{(13)}$ only the configuration
$(k,q,-q)$ is needed. The source $c^{(3)}$ is explicit
(\code{solution/results/third\_order\_sources.txt}) in the first-order closed forms and in
$K_2(k,q)$, $K_2(k,-q)$ and $K_2(q,-q)$. The last of these is the exact $K=0$ kernel, never a
limit. The third-order kernels are the formulas above applied to $c^{(3)}$, so they are double
integrals with closed-form inner integrals.

%======================================================================
\section{The observable through third order}\label{sec:obs}
%======================================================================

\begin{proposition}[Tetrad]\label{prop:frame}
Fermi--Walker transport along $u$ is parallel transport along $\partial_t$ in $\tilde g$, that
is, $\dd E_a{}^i/\dd t=-S^i{}_jE_a{}^j$. With $E(t_I)=\gamma(t_I)^{-1/2}$, the transported frame
satisfies $E=\gamma_o^{-1/2}R$ at observation, where $R=1+\eps^2\Omega/2+O(\eps^3)$ and $\Omega$
is antisymmetric. Consequently:
\begin{itemize}
\item $\Theta^{(1)}$ and $\Theta^{(2)}$ are those of the symmetric frame;
\item $\Theta^{(3)}$ differs only by an infinitesimal rotation of $\Theta^{(1)}$;
\item $C_\ell^{(22)}$ and $\operatorname{Re}C_\ell^{(13)}$ are individually unchanged.
\end{itemize}
\end{proposition}
\begin{proof}
At first order $E_1'=-\frac12\gamma_1'$, so $E_1=-\frac12\gamma_1$, which is the first-order
part of $\gamma^{-1/2}$. The rotation acts on $a^{(1)}_{\ell m}$ by an anti-Hermitian generator,
so $\operatorname{Re}\sum_ma^{(1)*}_{\ell m}(D(\Omega)a^{(1)})_{\ell m}=0$ in every realisation.
\end{proof}

Take the background ray $\bar X=\nh(t_o-t)$, $\bar e=-\nh$, and write $X=\bar X+\eps X_1+\eps^2X_2/2$.
All fields are evaluated on $(t,\bar X)$, and $\sigma(a,b)=\sigma_{ij}a^ib^j$. Because $\sigma$,
$S$ and $\Gamma$ vanish on the background, $X_3$ and $e_3$ never enter. The exact consequences of
Lemma~\ref{lem:redshift} are as follows.

\paragraph{Integrand.}
\begin{align*}
F_1&=\sigma_{1,nn},\qquad F_2=\sigma_{2,nn}+2X_1\!\cdot\!\partial\sigma_{1,nn}-4\sigma_1(n,e_1),\\
F_3&=\sigma_{3,nn}+3X_1\!\cdot\!\partial\sigma_{2,nn}-6\sigma_2(n,e_1)+3X_2\!\cdot\!\partial\sigma_{1,nn}+3(X_1\!\cdot\!\partial)^2\sigma_{1,nn}\\
&\quad-12(X_1\!\cdot\!\partial)\sigma_1(n,e_1)-6\sigma_1(n,e_2)+6\sigma_1(e_1,e_1).
\end{align*}

\paragraph{Ray.} The displacements are
\[
X_r=e_r(t_o)(t-t_o)+\int_t^{t_o}(t'-t)A_r(t')\dd t',
\]
with
\begin{align*}
A_1&=2S_1n-\Gamma_1(n,n)-n\sigma_{1,nn},\\
A_2&=[2S_2n-\Gamma_2(n,n)-n\sigma_{2,nn}]+2(X_1\!\cdot\!\partial)[2S_1n-\Gamma_1(n,n)-n\sigma_{1,nn}]\\
&\quad+2[-2S_1e_1+2\Gamma_1(n,e_1)+e_1\sigma_{1,nn}+2n\sigma_1(n,e_1)].
\end{align*}
The initial directions are the Taylor coefficients of $e(t_o)=-e^{-\zeta}\exp(-H/2)\,\nh$,
evaluated at $(t_o,0)$.

\paragraph{Emission and endpoint.} Let $\varphi_r$ denote $\nu$ along the ray, with
$\varphi_1=\nu_1$, $\varphi_2=\nu_2+2X_1\cdot\partial\nu_1$, and so on. Then
\begin{align*}
d_1&=-\bar t_\star\varphi_1,\qquad
d_2=\bar t_\star[u_1^2-\varphi_2-2\varphi_1'd_1],\qquad u_1=d_1/\bar t_\star,\\
I_1&=\int F_1,\qquad I_2=\int F_2-2F_1d_1,\qquad I_3=\int F_3-3F_1d_2-3F_2d_1-3F_1'd_1^2 .
\end{align*}

\paragraph{Clock.} With $\omega(t)=\nu(t,0)$ and $W_r(T)=\int_0^Tt(\omega_2+\omega_1^2)_r\dd t$,
the observation-time shifts are
\begin{align*}
D_1&=[\bar t_Id_{I1}-W_1(\bar t_o)+W_1(\bar t_I)]/\bar t_o,\\
D_2&=[\bar t_Id_{I2}+d_{I1}^2-W_2(\bar t_o)-2\bar t_o\omega_1(\bar t_o)D_1+W_2(\bar t_I)+2\bar t_I\omega_1(\bar t_I)d_{I1}-D_1^2]/\bar t_o .
\end{align*}
The coordinate-event map $\Theta_c$ is then composed with the random observation time:
\[
\Theta^{(2)}=\Theta_c^{(2)}+2D_1\partial\Theta_c^{(1)},\qquad
\Theta^{(3)}=\Theta_c^{(3)}+3D_2\partial\Theta_c^{(1)}+3D_1\partial\Theta_c^{(2)}+3D_1^2\partial^2\Theta_c^{(1)} .
\]

\paragraph{Normalisation.} Since $T\propto e^{-I}$, and $\sky{\cdot}$ is the random sky mean:
\begin{align*}
\Theta^{(1)}&=\sky{I_1}-I_1,\qquad
\Theta^{(2)}=I_1^2-I_2-2\sky{I_1}I_1+2\sky{I_1}^2-\sky{I_1^2}+\sky{I_2},\\
\Theta^{(3)}&=-I_3+3I_1I_2-I_1^3+3\sky{I_1}I_1^2-3\sky{I_1}I_2-3\sky{I_2}I_1+3\sky{I_1^2}I_1-6\sky{I_1}^2I_1\\
&\quad+\sky{I_3}-3\sky{I_1I_2}+\sky{I_1^3}+6\sky{I_1}^3-6\sky{I_1}\sky{I_1^2}+6\sky{I_1}\sky{I_2}.
\end{align*}
The dipole is removed last, by projection. It does not affect $C_{\ell\ge2}$, and no deboost is
implied.

\paragraph{Fourier dictionary.} For a single mode $a$, with phase $e^{i\lambda k\cdot\nh}$ and
$\lambda=t_o-t$:
\begin{itemize}
\item $\sigma_{a,nn}=\frac12Q_a'k^2(\mu_a^2-\frac13)$;
\item $\Gamma_a(n,n)=i[2(k\cdot n)n-k]+\frac i2[2(k\cdot n)H_an-kH_{a,nn}]$;
\item $e_a(t_o)=(1+\frac12H_a(t_o))\nh$;
\item $d_a=-\bar t_\star\nu_a(\bar t_\star)e^{i\lambda_\star k\cdot\nh}$.
\end{itemize}
The second-order shear is
$\sigma_{ab}=\frac12H_{ab}'+\frac14(H_aH_b+H_bH_a)'+H_a'+H_b'$, and
\begin{align*}
I_{ab}&=\int_{\bar t_\star}^{t_o}\!\dd t\,\big\{\sigma_{ab,nn}+i(k_b\!\cdot\!X_a)\sigma_{b,nn}+i(k_a\!\cdot\!X_b)\sigma_{a,nn}-2\sigma_b(n,e_a)-2\sigma_a(n,e_b)\big\}-[F_ad_b+F_bd_a]_{\bar t_\star},\\
\Theta_{ab}&=I_aI_b-I_{ab}-\sky{I_a}I_b-\sky{I_b}I_a+2\sky{I_a}\sky{I_b}-\sky{I_aI_b}+\sky{I_{ab}}+D_a\partial\Theta_b+D_b\partial\Theta_a .
\end{align*}
The third-order kernel $\Theta_{abc}$ is the multilinear coefficient of the formulas above.

\paragraph{Independent generator.} \code{observable\_generator.py} rebuilds $\Theta$ \emph{from
the exact definitions} by fixed-point iteration in the multilinear algebra: the ray ODE, the
emission root, the clock equation and the sky normalisation. It uses no hand-expanded formula,
and it passes the Ward identities and the superhorizon check of \S\ref{sec:gauge} and
\S\ref{sec:sh}.

%======================================================================
\section{Gauge invariance and Ward identities}\label{sec:gauge}
%======================================================================

Every coordinate and label is a physical scalar or frame: $t$, $X$, the surface $T=T_\star$,
the fluid line $X=0$, proper time from $T=T_I$, and $\gamma^{-1/2}$ in the labels. So $\Theta(\nh)$
is invariant under any diffeomorphism acting on fields and labels together. The only residual
freedom is the conformal group of the primordial chart (\S\ref{sec:ic}). Its dilation generator
must act as $\zeta\to\zeta+c\Leftrightarrow k\to e^{-c}k$. We verify this exactly:

\begin{itemize}
\item \emph{Fields:} $K_2(k,0)=-k\cdot\partial_kK_1$ and $K_3(k,0,0)=(k\cdot\partial_k)^2K_1$
  solve every equation in closed form. Perturbed predictions fail.
\item \emph{Observable:} $\Theta_{sL}=-\kappa\partial_\kappa\Theta_s$ and
  $\Theta_{sLL}=(\kappa\partial_\kappa)^2\Theta_s$ hold through $O(\kappa^4)$, where $\kappa$ scales
  the momenta. At this order the ray deflection, time delay, emission shift, clock, frame and
  normalisation all contribute.
\item At first order, the result equals the independent Newtonian-gauge computation
  (SW, Doppler and ISW) for every~$\ell$.
\end{itemize}

%======================================================================
\section{The spectrum}\label{sec:spec}
%======================================================================

\subsection{Gaussian reduction}

With $\mathcal K_r(k_1..k_r;\nh)$ the kernel of $\Theta^{(r)}$ (observer at the origin) and
$A^{(r)}_{\ell m}=\int\dd\Omega\,Y^*_{\ell m}\mathcal K_r$:
\begin{align*}
C^{(11)}_\ell&=\tfrac{1}{2\ell+1}\textstyle\sum_m\int_kP|A^{(1)}_{\ell m}|^2,\qquad
C^{(22)}_\ell=\tfrac{2}{2\ell+1}\sum_m\int_k\int_qP(k)P(q)|A^{(2)}_{\ell m}(k,q)|^2,\\
C^{(13)}_\ell&=\tfrac{3}{2\ell+1}\textstyle\sum_m\int_k\int_qP(k)P(q)A^{(1)}_{\ell m}(k)A^{(3)}_{\ell m}(k,q,-q)^* .
\end{align*}
The disconnected part of 22 is proportional to $\delta_{\ell0}$. In 13 all three pairings
survive; the internal $(q,-q)$ pairing is a pairing of the fourth moment and is not removed by
normal ordering. Odd orders vanish.

\subsection{Monochromatic spectrum}\label{sec:mono}

\emph{After} the complete observable is constructed, set
\[
\mathcal P(k)=\sigma^2\delta(\ln k/k_0),\qquad P(k)=\frac{2\pi^2\sigma^2}{k_0^2}\delta(k-k_0),
\qquad\ens{g^2}=\sigma^2\ \text{held fixed}.
\]
The background and the labels $T_\star$, $T_I$ and $\tau_o$ are also held fixed. Using joint
rotation invariance, with $\hat k=\hat z$ and $\hat q_\mu=(\sqrt{1-\mu^2},0,\mu)$:
\begin{align}
C^{(11)}_\ell&=4\pi\sigma^2T_\ell(k_0)^2,\\
C^{(22)}_\ell&=\frac{\sigma^4}{4\pi}\int_{-1}^1\!\dd\mu\!\int\!\dd\Omega\!\int\!\dd\Omega'\,P_\ell(\nh\cdot\nh')\,\mathcal K_2(\nh)\mathcal K_2(\nh')^*,\\
C^{(13)}_\ell&=\frac{3\sigma^4}{2}\,i^\ell T_\ell(k_0)\int_{-1}^1\!\dd\mu\Big[\int\!\dd\Omega\,P_\ell(\hat z\cdot\nh)\,\mathcal K_3(k_0\hat z,k_0\hat q_\mu,-k_0\hat q_\mu;\nh)\Big]^* .
\end{align}

\paragraph{All sky integrals in closed form.} Every term of $\mathcal K_r$ has the form
$c(\lambda,\ldots)\,\Pi(\nh)\,e^{i\nh\cdot V}$. Here $V=\sum_ak_a\lambda_a$ collects the
line-of-sight and nested times, and $\Pi$ is a polynomial. By Funk--Hecke and
$n_i\to-i\partial_{V_i}$:
\begin{align*}
\int\dd\Omega\,P_\ell(\hat z\cdot\nh)\,\Pi(\nh)e^{i\nh\cdot V}&=\Pi(-i\partial_V)\big[4\pi i^\ell j_\ell(|V|)P_\ell(\hat z\cdot\hat V)\big],\\
\int\!\!\int P_\ell(\nh\cdot\nh')\Pi(\nh)e^{i\nh\cdot V}\Pi'(\nh')e^{-i\nh'\cdot V'}&=\Pi(-i\partial_V)\Pi'(i\partial_{V'})\big[(4\pi)^2j_\ell(|V|)j_\ell(|V'|)P_\ell(\hat V\cdot\hat V')\big].
\end{align*}
The sky means in the normalisation terms are the $\ell=0$ case, $4\pi j_0$. The remaining
integrals are therefore:
\begin{itemize}
\item for $C^{(22)}$: $\mu$, plus per copy one line-of-sight time, one ray time and one
  field-kernel time;
\item for $C^{(13)}$: $\mu$ and up to five nested times.
\end{itemize}
All integrands are elementary or closed-form.

\begin{theorem}[Existence]\label{thm:exist}
For all finite $0<x_\star<x_o$, the integrals for $C_\ell^{(22)}$ and $C_\ell^{(13)}$ converge
absolutely. The restriction to the shell is well defined.
\end{theorem}
\begin{proof}
The domains are compact. The integrands are bounded for the following reasons:
\begin{itemize}
\item the field kernels are continuous in time and momenta, and finite on resonances, by the
  entire-function representation;
\item $c_M^{(2)}$ is analytic and vanishes at $K=0$, so $c_M=O(K)$, and the projector terms
  $(\hat K\hat K-\delta/3)h_S$ and $K\otimes w$ are bounded near $K\to0$. They are
  discontinuous only at $\mu=-1$, a null set;
\item $K_2(q,-q)$ in $C^{(13)}$ is the exact $K=0$ kernel;
\item the ray integrals are over finite intervals.
\end{itemize}
\end{proof}

%======================================================================
\section{Limits}\label{sec:limits}
%======================================================================

\subsection{Superhorizon limit, exact in \texorpdfstring{$\zeta$}{zeta}}\label{sec:sh}

\begin{theorem}\label{thm:sh}
Let all momenta scale as $\kappa\to0$ at fixed times. Then, to all orders in~$\eps$,
\[
\Theta(\nh)=-\tau^2\,e^{-2\zeta}\,\TF\big[\partial_i\partial_j\zeta-\partial_i\zeta\partial_j\zeta\big]_o\,\nh^i\nh^j+O(\kappa^3),
\qquad\tau^2=\tfrac16(\bar t_o^2-\bar t_\star^2).
\]
\end{theorem}
\begin{proof}
$\sigma=O(\kappa^2)$, while its spatial variation along the ray, the direction corrections and
the ray displacement all bring at least one more power of $\kappa$. The endpoint and clock
shifts are themselves $O(\kappa^2)$ and multiply $O(\kappa^2)$ quantities. Hence
$I=\frac12\big[(e^H)_{nn}\big]^{t_o}_{\bar t_\star}$ at $X=0$. Here $e^{-2\zeta}\sigma=\frac12(e^H)'$
and $e^{\zeta}e$ is $\exp(-H/2)\nh$ to this order. Inserting \eqref{eq:grad} gives $I$. Since
$I$ is a pure quadrupole, $\sky{I}=0$, and $\Theta=-I+O(\kappa^4)$.

The bound is uniform in the angles, because the remainders are bounded (projectors are
bounded), so the expansion may be taken inside the compact integrals.

The independent generator confirms the formula at first and second order, at $(k,q)$ with
general momenta. The corresponding third-order run at $(k,q,-q)$ was \emph{not completed}
(\S\ref{sec:ledger}). At third order the formula rests on two things: the proof above, which uses
only the gradient-expansion data \eqref{eq:grad} (exact in $\zeta$ and verified for arbitrary
$\zeta$), and the generator's third-order Ward identity W3.
\end{proof}

With $\Xi\equiv x_o^2-x_\star^2=k_0^2(\bar t_o^2-\bar t_\star^2)/3$, exact Wick moments of
$(g,\partial g,\partial\partial g)$ at a point give the leading $O(k_0^4)$ results \st{closed form}:
\[
C_2^{(11)}=\frac{4\pi}{225}\sigma^2\Xi^2,\qquad
C_2^{(22)}=\frac{272\pi}{675}\sigma^4\Xi^2,\qquad
C_2^{(13)}=\frac{16\pi}{75}\sigma^4\Xi^2,
\]
\[
C_2=\frac{4\pi}{225}\eps^2\sigma^2\Xi^2\Big[1+\frac{29}{3}\eps^2\sigma^2\Big]+O(\eps^6)+O(k_0^6),\qquad C_{\ell\ge3}=O(k_0^6).
\]
Two checks: $C_2^{(11)}$ agrees with the line-of-sight transfer function
$T_2\to-\frac{2}{15}k_0^2\tau^2$, and the direct $\eps$-expansion of $\ens{\sum_m|a_{2m}|^2}$
reproduces $\frac14C^{(22)}+\frac13C^{(13)}$. The fourth-order correction is a finite, positive
renormalisation of the quadrupole by the factor $1+\frac{29}{3}\eps^2\sigma^2$. It is
independent of $T_I$ at this order.

\subsection{Late time, first order}

\begin{proposition}
For $\ell\ge2$ and fixed $x_\star$, $T_\ell(x_\star,x_o)=O(\ln x_o/x_o)$ as $x_o\to\infty$. Hence
$C^{(11)}_\ell\to0$.
\end{proposition}
\begin{proof}
Write $T_\ell=-3\int_0^{x_o-x_\star}s'(x_o-u)w_\ell(\sqrt3u)\dd u$, with $w_\ell=j_\ell/3+j_\ell''$.
Two integrations by parts give $\int_0^\infty e^{iu}w_\ell(\sqrt3u)\dd u=0$ exactly for
$\ell\ge2$, because $j_\ell(0)=j_\ell'(0)=0$. This is the observer-local sound-cone
cancellation. Since $s'(x)=\sin x+O(1/x)$, the $O(1)$ oscillation therefore vanishes. The
remaining pieces are oscillatory with distinct frequencies $\sqrt3\pm1$ and amplitudes
$O(1/x_o)$, giving $O(\ln x_o/x_o)$ by integration by parts.
\end{proof}

At fourth order the late-time limit is non-uniform at the resonances $\mu=\pm1$ and
$\mu=-1/3$. Only the order ``angle first'' is admissible. Its evaluation is \st{open}.

\subsection{General spectra}

\begin{theorem}[Infrared]\label{thm:ir}
For spectra smooth on $(0,\infty)$, consider the $q\to0$ region. By the Ward identities,
$K_2(k,q)=-k\cdot\partial_kK_1(k)+O(q)$ and $K_3(k,q,-q)=(k\cdot\partial_k)^2K_1(k)+O(q^2)$.
Here the kernels are analytic in $q$ at fixed finite times, and the absence of an $O(q)$ term in
$K_3$ follows from its symmetry under $q\to-q$. The $O(q)$ cross term in $|K_2|^2$ is odd in
$\hat q$ and integrates to zero.

Both the $q$-long and the $k$-long regions of $C^{(22)}$ contribute. Let
$\sigma_L^2=\int_{q<\Lambda}P$ and $D=k\cdot\partial_k$. The infrared-sensitive part of
$\frac14C^{(22)}+\frac13\operatorname{Re}C^{(13)}$ is then
\[
\sigma_L^2\,\frac{1}{2\ell+1}\int_kP(k)\Big[\sum_m|DA^{(1)}_{\ell m}|^2+\operatorname{Re}\sum_mA^{(1)}_{\ell m}D^2A^{(1)*}_{\ell m}\Big]
=\frac{\sigma_L^2}{2}\,\frac{1}{2\ell+1}\int\frac{\dd^3k}{(2\pi)^3}\Big[\sum_m|A^{(1)}_{\ell m}|^2\Big](D+3)^2P(k).
\]
This vanishes identically for scale-invariant $P\propto k^{-3}$. Otherwise it is the physical
rescaling of the short-mode spectrum by the long-mode variance. The remainder is
$O(q^2)P(q)$, which is integrable for $\mathcal P(q)$ with $\int_0\mathcal P(q)q^2\,\dd\ln q<\infty$.
\end{theorem}

\paragraph{Ultraviolet \st{power counting only}.} At fixed finite times:
\begin{itemize}
\item the first-order kernels are bounded uniformly in $k$;
\item each order of the sources carries at most two spatial derivatives;
\item each ray displacement carries one.
\end{itemize}
Counting powers on this basis gives $|\mathcal K_2(k,q)|\lesssim q^2$ and
$|\mathcal K_3(k,q,-q)|\lesssim q^4$ as $q\to\infty$. That suggests absolute convergence of
$C^{(22)}$ and $C^{(13)}$ whenever $\int^\infty\mathcal P(q)q^4\,\dd\ln q<\infty$.

This is \emph{not} a proof. A rigorous bound needs uniform estimates of the lapse Green function
near $s\to0$ (where $|G_N|$ grows like $1/(K^3ts)$) against sources whose small-$s$ Taylor
coefficients grow with $q$. The true large-$q$ behaviour for a scale-invariant spectrum,
including possible cancellations, is \st{open}.

%======================================================================
\section{Verification ledger}\label{sec:ledger}
%======================================================================
All checks are exact symbolic computations (SymPy~1.14). None uses floating point or
quadrature as evidence. ``Generic 2-jets'' means that every function value and derivative is
an independent symbol, so a zero residual is a proof.

\begin{longtable}{p{6.2cm}p{8.8cm}}
\toprule
Check & Script and scope\\\midrule\endhead
Lemmas L1--L5; generic field equations & \code{verification/check\_exact\_lemmas.py}, \code{check\_fw\_conformal.py}, \code{check\_field\_equations\_generic.py} (generic 2-jets)\\
Mixed trace-free form (T) & \code{solution/check\_evolution\_form.py} (generic 2-jets)\\
First order, two gauges, all $\ell$ & \code{verification/check\_first\_order.py}; \code{solution/first\_order\_engine\_check.py}\\
Closed form $T_2$, $T_3$ & \code{verification/transfer\_closed\_form.py}\\
Green functions, $K\to0$ & \code{solution/check\_green.py}\\
2nd-order sources and operators & \code{solution/second\_order.py}\\
2nd order, all 10 equations & \code{solution/second\_order\_consistency.py}: exact Taylor series through $O(t^8)$, general $K$ and $K=0$\\
2nd-order closed forms & \code{solution/closed\_second\_order.py}: every antiderivative by exact differentiation; (N) by closed-form substitution (\code{--check N}); (M) and (T) via the theorem and the series test\\
3rd-order sources and operators & \code{solution/third\_order.py}\\
3rd order, all 10 equations & \code{solution/fast\_series.py 8}: the pair kernels, including the exact $K=0$ pair, then the triple at $(k,q,-q)$ with general momenta. The \emph{full nonlinear} equations are recomputed with the solved kernels and vanish through $O(t^6)$. The arithmetic is exact (sparse polynomials over $\mathbb Q(i)$ with formal inverse momenta; denominators are cleared exactly at the zero test). Negative controls (a wrong $c_3$, a corrupted $K=0$ pair kernel) fail at the expected orders. The run takes about 70~s. \code{third\_order\_consistency.py} is the slower SymPy reference implementation of the same test.\\
Dilation Ward identities (fields) & \code{solution/check\_ward\_dilation.py}, \code{check\_ward\_dilation3.py} (closed form; negative controls)\\
Observable expansion formulas & \code{solution/check\_observable\_expansion.py} (arbitrary fields, general jets)\\
Observable Ward identities & \code{solution/observable\_generator.py W2|W3} through $O(\kappa^4)$\\
Superhorizon observable & \code{solution/observable\_generator.py SH}: orders 1 and 2, general $(k,q)$. \emph{Not completed:} \code{SH3} (third order at $(k,q,-q)$). Its fields are built by the fast engine in about 12~s, but the generator's SymPy assembly of the observable did not finish within 50~minutes. Porting the generator to the fast arithmetic would remove this; it was not done.\\
Superhorizon $C_2$ & \code{solution/superhorizon.py}\\
Normalisation; Wick and Fourier measures & \code{verification/normalization\_expansion.py}; \code{validation/check\_algebra.py}\\
\bottomrule
\end{longtable}

%======================================================================
\section{What remains unevaluated}\label{sec:open}
%======================================================================
\begin{enumerate}
\item For general $k_0$, $C_\ell^{(22)}$ and $C_\ell^{(13)}$ are given as exact absolutely
  convergent integrals: over $\mu$ and the nested times, with closed-form integrands, and with
  all sky integrals evaluated. We do not expect, and did not find, a closed form. They are
  \emph{not} evaluated beyond the superhorizon limit.
\item The late-time limit at fourth order (non-uniform at the resonances) is not evaluated.
\item Ultraviolet convergence for spectra without compact support is established by power
  counting only, not by a rigorous bound.
\item The infrared statement (Thm.~\ref{thm:ir}) assumes that the kernels are analytic in $q$ at
  $q=0$ at fixed finite times. This holds for the integral representations by construction, but
  the uniformity of the expansion in $k$ over unbounded spectra is not proved.
\item The perturbative domain is treated order by order: a single emission intersection and no
  caustics. There is no statement about convergence in~$\eps$.
\end{enumerate}

\end{document}
